\section{RCOUNT Inputs And Report Contract}

RCOUNT can feed this method with normalized summaries and large count tables:
contest totals by reporting unit, jurisdiction, candidate, batch, or other
publicly reproducible grouping. The method must bind every report to the source
package hashes so that a second analyst using the same RCOUNT package and the
same registered method recomputes the same extracted digits and scores
\citep{rcount-overview-2026}.

W.04 instantiates the W.01 transcript fields as follows.

\begin{center}
\small
\begin{tabularx}{\textwidth}{lX}
\toprule
Field & W.04 instantiation \\
\midrule
\texttt{method\_id} & Registered digit-distribution screen. \\
\texttt{feature\_set} & Count column, digit position, exclusions, pooling rule,
and any minimum-count threshold. \\
\texttt{baseline\_population} & Contest, jurisdiction set, election, time window,
and comparable count table used for the reference. \\
\texttt{model\_id} & \texttt{digit-test/benford}, with subtype for first digit,
second digit, last digit, or declared distance family. \\
\texttt{unit\_id} & Contest, jurisdiction, reporting unit, batch, or aggregate
count table being scored. \\
\texttt{score} & Test statistic or distance from the declared reference. \\
\texttt{p\_value\_or\_rank} & Optional p-value, percentile, or queue rank, reported
only with its assumptions. \\
\texttt{investigation\_note} & Required weak-signal caveat: digit tests are
low-weight screens, deviations are expected under bounded or similar precinct
sizes, and the result is not evidence of fraud. \\
\bottomrule
\end{tabularx}
\end{center}

The source hashes, feature definition, baseline population, and model identity
are part of the evidence. Changing any of them changes the question. The
investigation note is not optional boilerplate; it is the field that prevents a
large digit statistic from being circulated as a certification or fraud claim.
